\AtBeginDocument{%
  \providecommand\BibTeX{{%
    Bib\TeX}}}

% K-LD macros (also \provided in the body sections; defined here so they are
% available in the title/abstract, which are typeset before the body is input).
\providecommand{\kld}{\textsc{K-LD}}
\providecommand{\kst}{\textsc{K-ST}}

\setcopyright{acmlicensed}
\copyrightyear{2026}
\acmYear{2026}
\acmDOI{XXXXXXX.XXXXXXX}

\acmJournal{TCPS}
\acmVolume{1}
\acmNumber{1}
\acmArticle{1}
\acmMonth{6}

\title[\kld{}: Executable Semantics for IEC~61131-3 Ladder Diagram]{\kld{}: An Executable Formal Semantics of IEC~61131-3\\Ladder Diagram for Validating Verifier Translations}

\renewcommand{\shortauthors}{Anonymous}

\begin{abstract}
Automated verifiers for \gls{iec}~61131-3 \gls{ld} establish safety by translating
a diagram into the input model of a model checker---the \gls{esbmc}-\gls{plc} line,
for example, lowers \gls{ld} into a GOTO \gls{ir} and discharges the result with
\gls{smt}-based bounded model checking. Every such tool rests on an unstated
assumption: that this \emph{front-end} faithfully reflects the standard. The
translation is nonetheless unverified code, and when it disagrees with
\gls{iec}~61131-3 the verifier silently returns the wrong answer---missing real
violations or raising false alarms---with no independent account of what the
diagram \emph{means} to catch the error. We close this gap with \textbf{\kld{}}, an
\emph{executable} formal semantics of \gls{iec}~61131-3 \gls{ld} in the \textsf{K}
framework. Extending \kst{} to the graphical language, \kld{} models contacts,
energise/latch/unlatch coils, the retentive scan cycle, the timers
(\code{TON}/\code{TOF}/\code{TP}), counters (\code{CTU}/\code{CTD}), and edge
function blocks; from one definition \textsf{K} yields both an interpreter and a
deductive verifier. We validate \kld{} scan-for-scan against the OpenPLC/\gls{matiec}
reference implementations, then use it as an \emph{independent reference oracle} to
differentially validate the \gls{esbmc}-\gls{plc} \gls{ld}$\rightarrow$GOTO
translation over a benchmark suite. \kld{} agrees with \gls{esbmc} on the large
majority of programs and reproduces every injected violation with a concrete
witness; every \emph{disagreement} is a genuine \gls{esbmc} defect that \kld{}
exposes and OpenPLC confirms, falling into two opposite failure modes of its
incomplete timer translation---an unsound \emph{skip} that certifies a violating
program safe, and an imprecise \emph{havoc} that manufactures impossible
counterexamples. Finally, we \emph{machine-check} in \code{kprove} that \kld{}'s
rules implement the standard's input/output relation for the combinational and
latch fragment, raising the correctness argument there from empirical to formally
proven. \kld{} is a reusable, standard-faithful oracle for scrutinising \emph{any}
\gls{ld} verifier's front-end.
\end{abstract}

\begin{CCSXML}
<ccs2012>
   <concept>
       <concept_id>10011007.10010940.10010992.10010998.10010999</concept_id>
       <concept_desc>Software and its engineering~Software verification</concept_desc>
       <concept_significance>500</concept_significance>
       </concept>
   <concept>
       <concept_id>10010583.10010717.10010721.10003791</concept_id>
       <concept_desc>Hardware~Model checking</concept_desc>
       <concept_significance>500</concept_significance>
       </concept>
   <concept>
       <concept_id>10010520.10010553</concept_id>
       <concept_desc>Computer systems organization~Embedded and cyber-physical systems</concept_desc>
       <concept_significance>500</concept_significance>
       </concept>
   <concept>
       <concept_id>10011007.10010940.10011003.10011114</concept_id>
       <concept_desc>Software and its engineering~Software safety</concept_desc>
       <concept_significance>500</concept_significance>
       </concept>
   <concept>
       <concept_id>10011007.10011006.10011050.10011017</concept_id>
       <concept_desc>Software and its engineering~Domain specific languages</concept_desc>
       <concept_significance>500</concept_significance>
       </concept>
   <concept>
       <concept_id>10003752.10010124.10010138.10010142</concept_id>
       <concept_desc>Theory of computation~Program verification</concept_desc>
       <concept_significance>300</concept_significance>
       </concept>
   <concept>
       <concept_id>10003752.10003790.10002990</concept_id>
       <concept_desc>Theory of computation~Logic and verification</concept_desc>
       <concept_significance>300</concept_significance>
       </concept>
   <concept>
       <concept_id>10011007.10011074.10011099.10011102.10011103</concept_id>
       <concept_desc>Software and its engineering~Software testing and debugging</concept_desc>
       <concept_significance>300</concept_significance>
       </concept>
 </ccs2012>
\end{CCSXML}

\ccsdesc[500]{Software and its engineering~Software verification}
\ccsdesc[500]{Hardware~Model checking}
\ccsdesc[500]{Computer systems organization~Embedded and cyber-physical systems}
\ccsdesc[500]{Software and its engineering~Software safety}
\ccsdesc[500]{Software and its engineering~Domain specific languages}
\ccsdesc[300]{Theory of computation~Program verification}
\ccsdesc[300]{Theory of computation~Logic and verification}
\ccsdesc[300]{Software and its engineering~Software testing and debugging}

\keywords{executable formal semantics, \textsf{K} framework, \acrshort{iec}~61131-3, Ladder Diagram, \acrshort{plc}, reference oracle, differential testing, translation validation, scan cycle, function blocks, \acrshort{esbmc}, \acrshort{bmc}, model checking, OpenPLC, PLCopen XML}
